\documentclass[11pt]{article}
\usepackage[a4paper,margin=1in]{geometry}
\usepackage{amsmath,amssymb,amsthm}
\usepackage{mathtools}
\usepackage{booktabs}
\usepackage{graphicx}
\usepackage{xcolor}
\usepackage{hyperref}
\hypersetup{colorlinks=true,linkcolor=blue,urlcolor=blue,citecolor=blue}
\newcommand{\rt}[1]{\texttt{#1}}
\newcommand{\fn}[1]{\texttt{#1}}
\newcommand{\K}{\mathcal{K}}
\newcommand{\ifgraphic}[2]{\IfFileExists{#1}{\includegraphics[width=#2]{#1}}%
  {\fbox{\parbox{#2}{\centering\small figure \texttt{\detokenize{#1}} pending}}}}
\newcommand{\iftable}[1]{\IfFileExists{#1}{\input{#1}}%
  {\fbox{\parbox{0.9\textwidth}{\centering\small table
  \texttt{\detokenize{#1}} pending}}}}

\theoremstyle{plain}
\newtheorem{theorem}{Theorem}
\newtheorem{proposition}{Proposition}
\newtheorem{corollary}{Corollary}
\theoremstyle{remark}
\newtheorem{remark}{Remark}

\title{\textbf{Report R30 --- the sector-size distribution of the
one-Hadamard rules $W_{156}/W_{198}$ and $W_{108}/W_{201}$}\\
\large exact at finite $N$, log-normal at large $N$, at \rt{obc0} and \rt{pbc}}
\author{Tier 1 analysis}
\date{2026-08-28}

\begin{document}
\maketitle

\begin{abstract}
R9 counted the Krylov sectors of the four one-Hadamard rules and measured the
largest one; R18 showed both follow from a frozen wall word. This report gives
the \emph{whole} distribution: how many sectors there are of each size, at every
$N$, at both boundary conditions, and what that distribution looks like as
$N\to\infty$. Everything reduces to one $2\times2$ matrix. Because the wall word
has length two, its indicator matrix $M_1=|a\rangle\langle b|$ has \textbf{rank
one}, so a wall cuts the transfer product and the sector size factorises over
the gaps between consecutive walls,
$|\K| = L(g_0)K(g_1)\cdots K(g_{k-1})R(g_k)$, with all three kernels corners of
powers of the single matrix $M_0 = J - |a\rangle\langle b|$. The two families are
then exactly the two \textbf{Jordan types} of $M_0$: the words \rt{00} and
\rt{11} are diagonal, $M_0$ is a Fibonacci matrix and $K(g)=F(g-1)$ grows
exponentially; the words \rt{01} and \rt{10} are off-diagonal, $M_0$ is
unipotent and $K(g)=g$ grows linearly. From that one bit follow the sector
counts ($\varphi$ against $\rho^2$), the largest sector ($\varphi^N$ against
$4^{N/5}$), the frozen-sector counts (exponentially many against linearly many),
and the set of sizes that occur at all --- every integer for
$W_{156}/W_{198}$, but for $W_{108}/W_{201}$ only products of Fibonacci numbers,
joined on the ring by the one wall-free sector of Lucas size $L(N)$, so no sector
of those rules ever has size $14$, $17$, $19$, $22$, $23$, $28$ or $31$.
We give the number of sectors of a given size $\sigma$ in closed form: at
\rt{obc0} in the unipotent family
$n_N(\sigma)\simeq P(\sigma)(N/2)^{\Omega(\sigma)+1}/(\Omega(\sigma)+1)!$, a
polynomial whose degree is the number of prime factors of $\sigma$; on the ring
the same law acquires a strict parity selection rule; in the Fibonacci family
$n_N(\sigma)\simeq A(\sigma)N^{m(\sigma)}\varphi^N$. For the distribution as a
whole, $\lambda_s$, defined by the single scalar equation $B_s(1/\lambda_s)=1$
with $B_s(x)=\sum_g K(g)^sx^{g+1}$, is the free energy: $\lambda_0$ is the
sector-count base, $\lambda_1=2$ identically, and $\tau=\ln\lambda$ gives a
\textbf{log-normal} law $|\K| = \exp(\alpha_{\rm typ}N + \sigma\sqrt{N}\,\xi)$
together with its full large-deviation spectrum. $\lambda_s$ is
boundary-independent; the boundary contributes only an $O(1)$ offset, which
vanishes identically on the ring and is rule-dependent on the chain --- negative
for $W_{156}/W_{198}$ and $W_{108}$, but \emph{positive} for $W_{201}$.
Every step is checked against exhaustive enumeration of the $2^N$ flip graph,
and independently against the transfer-matrix construction of Tier-2 R-T20.
\end{abstract}

\tableofcontents

\section{The question, and what was already settled}

The four rules carry a single Hadamard that fires on one neighbour pattern:
\begin{center}
\begin{tabular}{lllll}
\toprule
rule & tuple & fires when $(x_{i-1},x_{i+1})=$ & wall word & family \\
\midrule
$W_{201}$ & \rt{VIII} & $(0,0)$ & \rt{11} & the PXP constraint \\
$W_{108}$ & \rt{IIIV} & $(1,1)$ & \rt{00} & its spin-flip image \\
$W_{156}$ & \rt{IIVI} & $(1,0)$ & \rt{01} & chiral \\
$W_{198}$ & \rt{IVII} & $(0,1)$ & \rt{10} & its reflection \\
\bottomrule
\end{tabular}
\end{center}

\noindent All four are unitary, so by R17 their terminal SCCs are their Krylov
sectors, and by the flip-reduction theorem (\fn{graph/flip\_graph.py}) the
sector partition is the connected-component partition of the single-flip graph
$x \sim x\oplus 2^i$ whenever the gate fires at $i$ in $x$. R9 fitted the
sector-count bases $\rho^2=1.754878$ and $\varphi$; R18 derived them from a
frozen wall word and showed that the conserved domain-wall charges are functions
of the wall set rather than a second ingredient. R21 classified all four as
strongly fragmented.

What none of those settled is the \emph{distribution}: how the $2^N$ basis states
are apportioned among the sectors. That is the object below. The user's question
was asked for both boundary conventions, and the answer differs between them in
ways that turn out to be sharp and worth stating separately, so \rt{obc0} and
\rt{pbc} are carried side by side throughout.

\section{The label is the sector}

Write the padded chain $y_0=0,\ y_1\ldots y_N = x,\ y_{N+1}=0$ at \rt{obc0}, and
$y_i = x_{i\bmod N}$ at \rt{pbc}. The \textbf{label} of a state is the set of
positions at which the rule's wall word occurs.

\begin{theorem}[completeness of the label]
\label{thm:label}
For each of the four rules, two states lie in the same Krylov sector if and only
if they have the same label --- with one exception, at \rt{pbc} for
$W_{156}/W_{198}$, where $0^N$ and $1^N$ share the empty label but are two
distinct frozen sectors.
\end{theorem}

\begin{proof}[Proof]
\emph{Constancy.} No move creates or destroys an occurrence. Take $W_{201}$,
word \rt{11}: a site flips only when both neighbours are $0$. If $y_i=y_{i+1}=1$
then site $i$ cannot flip (its right neighbour is $1$) and neither can site
$i+1$, so an existing wall is frozen. Conversely a new \rt{11} at $j$ would need
some site to become $1$ next to a $1$; but flipping $x_j$ up requires
$x_{j-1}=x_{j+1}=0$, so no wall is created at $j-1$ or $j$. For $W_{156}$, word
\rt{01}: a wall at $i$ means $y_i=0,y_{i+1}=1$; site $i$ can flip only if
$(y_{i-1},y_{i+1})=(1,0)$, contradicted by $y_{i+1}=1$, and site $i+1$ only if
$(y_i,y_{i+2})=(1,0)$, contradicted by $y_i=0$. Creation is excluded the same
way. $W_{108}$ and $W_{198}$ are the images of these two under the spin flip and
the reflection of the local rule, and the same two lines apply verbatim.

\emph{Transitivity.} Fix a label. For $W_{156}/W_{198}$ the label is the set of
ascents; ascents and descents alternate, so the free data are the descent
positions, one in each open interval between consecutive ascents. A descent at
$q$ can move to $q+1$ (flip $y_{q+1}$ up, which needs $(y_q,y_{q+2})=(1,0)$, true
inside the interval) and to $q-1$ (flip $y_q$ down, which needs
$(y_{q-1},y_{q+1})=(1,0)$, true inside the interval), and cannot leave it.
Distinct descents never interact because the frozen ascents separate them, so the
sector is a product of intervals and single flips connect it.
For $W_{201}$ the label fixes the maximal runs of $1$s of length $\ge 2$; the
sites adjacent to such a run are forced to $0$ and cannot flip, and the remaining
free stretches carry exactly the strings with no \rt{11}. In such a string every
$1$ has both neighbours $0$, so it may be flipped down; removing $1$s one at a
time reaches the all-zero configuration of that stretch without ever creating a
new run, so each free stretch is connected. $W_{108}$ is the same with $0$ and
$1$ exchanged.

\emph{The exception.} At \rt{pbc} the wall-free label of an ascent word is
realised by $\operatorname{tr}M_0^N=2$ states, $0^N$ and $1^N$, both frozen. Every
other wall-free fibre in this paper is connected.
\end{proof}

\noindent Theorem~\ref{thm:label} upgrades R18's exhaustive check to $N=14$ to a
proof. It is independently re-verified here by \fn{sector\_is\_wall\_set}, which
computes the label of all $2^N$ states and checks that it is constant on each
flip-graph component and injective across components, for all four rules at both
boundary conditions and $N=3\ldots18$: constant everywhere, zero collisions
except the single expected one for $W_{156}/W_{198}$ at \rt{pbc}.

\section{Rank one, and the renewal factorisation}

The wall word reads the pair $(y_i,y_{i+1})$, so the label letter at $i$ is a
function of one \emph{bond}. Hence the preimage of a label word $w$ is the set
of paths of a two-node graph with position-dependent edge sets,
\[
M_w[u][v] = 1 \iff \bigl[\,(u,v)=\text{word}\,\bigr]=w ,
\]
and, writing the word as $(a,b)$,
\begin{equation}
\label{eq:mats}
M_1 = |a\rangle\langle b| , \qquad M_0 = J - |a\rangle\langle b| ,
\qquad M_0+M_1 = J ,
\end{equation}
$J$ the all-ones matrix. The sector size is then a matrix element,
\[
|\K(w)| \;=\; \langle 0| M_{w_0}M_{w_1}\cdots M_{w_N} |0\rangle \quad(\rt{obc0}),
\qquad
|\K(w)| \;=\; \operatorname{tr} M_{w_0}\cdots M_{w_{N-1}} \quad(\rt{pbc}).
\]

\noindent $M_1$ is a \emph{matrix unit}, so it has rank one, and that is the
whole mechanism: every wall cuts the product.

\begin{proposition}[renewal factorisation]
\label{prop:renewal}
Let the walls sit at $p_1<\cdots<p_k$ and let $g_j$ be the number of non-wall
letters in the $j$-th gap, so that $g=0$ means two adjacent walls. Then
\begin{align}
|\K| &= L(g_0)\,K(g_1)\cdots K(g_{k-1})\,R(g_k) &&(\rt{obc0},\ k\ge1),\\
|\K| &= K(g_1)\cdots K(g_k) \quad\text{cyclically} &&(\rt{pbc},\ k\ge1),
\end{align}
with
\[
K(g)=(M_0^{\,g})_{b,a},\qquad L(g)=(M_0^{\,g})_{0,a},\qquad
R(g)=(M_0^{\,g})_{b,0},
\]
and the wall-free label has size $T(m)=(M_0^{\,m})_{0,0}$ at \rt{obc0} and
$\operatorname{tr}M_0^{\,m}$ at \rt{pbc}, with $m$ the number of label letters,
$N+1$ and $N$ respectively.
\end{proposition}

\noindent This is one formula for all four rules and both boundary conditions.
Everything else in this report is a statement about powers of a single
$2\times 2$ matrix.

\subsection{The two Jordan types}

\begin{center}
\small
\iftable{tab_r30_kernels.tex}
\end{center}

\noindent The word is \textbf{diagonal} ($a=b$) exactly for $W_{108}$ and
$W_{201}$ and \textbf{off-diagonal} for $W_{156}$ and $W_{198}$, and that one bit
is the fork:
\[
\begin{aligned}
\text{diagonal word:}\quad & M_0 \sim \begin{psmallmatrix}1&1\\1&0\end{psmallmatrix},
&& \text{eigenvalues } \varphi,\,-1/\varphi, && K(g)=F(g-1) && \text{HYPERBOLIC},\\
\text{off-diagonal word:}\quad & M_0 \sim \begin{psmallmatrix}1&0\\1&1\end{psmallmatrix},
&& \text{eigenvalue } 1 \text{ (defective)}, && K(g)=g && \text{PARABOLIC}.
\end{aligned}
\]

\begin{remark}
This is also why $W_{156}$ and $W_{198}$ have identical size multisets at every
$N$ and both boundary conditions, and why $W_{108}$ and $W_{201}$ do at
\rt{pbc}: the two $M_0$ of a pair are related by transposition or by conjugation
with the swap, which moves $(a,b)$ with them, so $K$ is the same corner. At
\rt{obc0} the boundary vector $|0\rangle$ is not swap-invariant, which is exactly
why the spin flip survives on the ring and dies on the chain. Nothing in this
report was obtained by \emph{applying} a symmetry: all four rules were computed
from their own wall words, and the coincidences are results.
\end{remark}

\subsection{Six consequences, read straight off the kernel}

\begin{enumerate}
\item \textbf{The largest sector is the wall-free label} in the hyperbolic
family, because a wall costs a letter and cuts the Fibonacci growth. Its size is
the corresponding corner of $M_0^{\,m}$:
\[
D_{\max}(W_{201},\rt{obc0}) = F(N+2), \quad
D_{\max}(W_{108},\rt{obc0}) = F(N), \quad
D_{\max}(W_{108/201},\rt{pbc}) = L(N),
\]
$L$ the Lucas numbers. The two Fibonacci indices between the two chain values are
the whole \rt{obc0} asymmetry: the frozen $0$-pad is a free \emph{collar} for the
word \rt{11} (the site next to a block is forced to $0$ anyway, so the pad costs
nothing) but a \emph{block letter} for the word \rt{00} (the wall-free stretch
must begin and end with a $1$, and the two pads are $0$s). All three grow like
$\varphi^N$.

\item \textbf{In the parabolic family a wall-free stretch is worth only its
length}, so the biggest sector is the best product of gap lengths at a cost of
one extra letter per gap:
$D_{\max} = \max \prod_j g_j$ subject to $\sum_j (g_j+1) \le m$. Since
$4^{1/5}=1.319508 > 3^{1/4}=1.316074$, the optimum uses gaps of $4$ and
$D_{\max}\sim 4^{N/5}$. Tier-2 R-T20 gives the exact residue form
$G(B)=3^k4^{(B-4k)/5}$ with $k=4B \bmod 5$, valid for $B\ge12$, and
$G(B+5)=4G(B)$ exactly. The series is
$2,3,4,5,6,9,12,16,20,27,36,48,64,81$ at $N=3\ldots16$ on the ring, and the same
one site later on the chain.

\item \textbf{Frozen sectors.} A sector has size $1$ iff every renewal factor is
$1$. In the hyperbolic family $K(g)=1$ for $g\in\{0,2,3\}$ (since
$F(-1)=F(1)=F(2)=1$), and $x+x^3+x^4=1$ at $x=1/\varphi$, so there are
$\Theta(\varphi^N)$ frozen sectors. In the parabolic family only $g=1$ gives
$K=1$, which pins the walls to a rigid alternating pattern and leaves
\emph{linearly} many: exactly $\lfloor (N-1)/2\rfloor+2$ at \rt{obc0}, and
exactly $4$ at even $N$ and $2$ at odd $N$ on the ring (the two of them being the
ring defect of Theorem~\ref{thm:label}).

\item \textbf{The sector count.} At $s=0$ the renewal sum counts realisable
labels. Parabolic: $K(g)>0$ iff $g\ge1$, i.e.\ no two adjacent walls, so the
labels are the independent sets of a path or a cycle, giving $F(N+2)$ at
\rt{obc0} and $L(N)+1$ at \rt{pbc} (the $+1$ is the ring defect). Hyperbolic:
$K(g)=0$ only at $g=1$, so the labels are the binary words avoiding \rt{101},
whose generating function has denominator $1-2x+x^2-x^3$ --- R18's recurrence
$a(N)=2a(N-1)-a(N-2)+a(N-3)$, dominant root $\rho^2$, at \emph{both} boundary
conditions.

\item \textbf{The ring defect is the trace of a unipotent matrix.} The wall-free
fibre at \rt{pbc} has $\operatorname{tr}M_0^{\,N}$ elements, which is $2$ when
$M_0$ is unipotent and $L(N)$ when it is the Fibonacci matrix. So the anomaly R18
recorded as an exception is not a special case to remember.

\item \textbf{Which sizes occur at all.} The achievable sizes are the
multiplicative semigroup generated by the kernel values, \emph{together with the
wall-free label}, which is the one term that is not a product of kernels.
Parabolic: $K(g)=g$ takes every value, so every integer up to $N$ is a sector
size (verified: no integer in $1\ldots N$ is missing at $N=18,22,24$).
Hyperbolic: the kernel values are the \emph{Fibonacci numbers}, so
\[
\{\text{sizes}\} =
\begin{cases}
\{\text{products of } F(j),\ j\ge3\} & (\rt{obc0}),\\[2pt]
\{\text{products of } F(j),\ j\ge3\}\ \cup\ \{L(N)\} & (\rt{pbc}).
\end{cases}
\]
The extra ring term is the wall-free sector, which is a \emph{trace} rather than
a corner of $M_0^{\,N}$ and is therefore Lucas, not Fibonacci. Since Lucas
numbers are generally not products of Fibonacci numbers this is an infinite
family of extra sizes, not a handful of exceptions: $7=L(4)$, $11=L(5)$,
$29=L(7)$, $76=L(9)=4\cdot19$, $123=L(10)=3\cdot41$ and $199=L(11)$ all occur,
exactly once each, at that one $N$. What is excluded at \emph{both} boundaries is
an integer that is neither a Fibonacci product nor a Lucas number: $14$, $17$,
$19$, $22$, $23$, $28$, $31$, $33$, $35$, $37$, $38,\ldots$ --- none of these is
a sector size of $W_{108}$ or $W_{201}$ at any $N$.
\end{enumerate}

\section{Generating functions and the exact census}

Give a wall together with the gap after it the weight $x^{g+1}$, one variable per
label letter, and set
\[
B_s(x)=\sum_{g\ge0}K(g)^s x^{g+1},\quad
A_s(x)=\sum_{g\ge0}L(g)^s x^{g},\quad
C_s(x)=\sum_{g\ge0}R(g)^s x^{g+1},\quad
E_s(x)=\sum_{m\ge0}T(m)^s x^{m}.
\]
With $Z_N(s)=\sum_{\text{sectors}}|\K|^s$,
\begin{equation}
\label{eq:gf}
\sum_N x^{N+1}Z^{\rt{obc0}}_N(s) = E_s + \frac{A_sC_s}{1-B_s},
\qquad
\sum_N x^{N}Z^{\rt{pbc}}_N(s) = \frac{x B_s'(x)}{1-B_s} + \sum_N (\operatorname{tr}M_0^N)^s x^N .
\end{equation}
The check that these are the right objects is $s=1$: since $M_0+M_1=J$ by
\eqref{eq:mats}, the $s=1$ transfer operator is $J$ itself and $Z_N(1)=2^N$ is a
theorem rather than a test. In the renewal language the same statement is
$B_1(x)=x^2/(1-x)^2$ for the parabolic family, whose root is $x=1/2$ exactly.

\begin{center}
\small
\textbf{\rt{obc0}} --- $n_{\rm sectors}$ / $D_{\max}$ / $n_{\rm frozen}$\\[2pt]
\resizebox{\textwidth}{!}{\iftable{tab_r30_census_obc0.tex}}\\[8pt]
\textbf{\rt{pbc}} --- $n_{\rm sectors}$ / $D_{\max}$ / $n_{\rm frozen}$\\[2pt]
\resizebox{\textwidth}{!}{\iftable{tab_r30_census_pbc.tex}}
\end{center}

\noindent The exact size multiset is computed at any $N$ by
\fn{sector\_sizes.hist}, and agrees term by term with exhaustive union-find on
the $2^N$ flip graph for all four rules, both boundary conditions and
$N=3\ldots16$ --- $112$ cases, zero mismatches, full multisets and not merely
counts and maxima.

\section{$n_N(\sigma)$: how many sectors have a given size}

Proposition~\ref{prop:renewal} says a sector of size $\sigma$ is an \emph{ordered
factorisation} of $\sigma$ into gap kernels, padded out with kernel-$1$ gaps to
fill the chain. That gives $n_N(\sigma)$ exactly, and it is cheap: the dynamic
program only ever needs to carry a \emph{divisor} of $\sigma$, so it costs
$O(N\,d(\sigma))$ and runs at $N=10^4$ without effort, where the full histogram
--- which must carry every partial product --- is already impossible past
$N\approx100$.

\subsection{The parabolic family: a polynomial whose degree is number theory}

Here the kernel is $K(g)=g$, so a factor $f$ costs $f+1$ letters and the padding
factor $1$ costs $2$. The longest factorisations are into \emph{primes}, and they
dominate. Writing $\Omega(\sigma)$ for the number of prime factors with
multiplicity and $P(\sigma)=\Omega!/\prod_i a_i!$ for the number of orderings of
the prime multiset,
\begin{equation}
\boxed{\;
n_N(\sigma) \;\simeq\; P(\sigma)\,\frac{(N/2)^{\Omega(\sigma)+1}}{(\Omega(\sigma)+1)!}
\;}
\qquad (\rt{obc0},\ \sigma \text{ fixed},\ N\to\infty).
\end{equation}
So the sector census reads off the prime factorisation of the size. $n_N(4)$ and
$n_N(6)$ both grow like $N^3$ because $\Omega=2$ for both, but $n_N(6)$ is
asymptotically \emph{twice} $n_N(4)$ because $6=2\cdot3$ can be ordered two ways
and $4=2\cdot2$ only one; a prime size grows only like $N^2$. Measured
$n_N(\sigma)$ divided by this prediction, at \rt{obc0}:
\begin{center}
\small
\begin{tabular}{lrrrrrr}
\toprule
$\sigma$ & $\Omega$ & $P$ & $N=50$ & $N=200$ & $N=800$ & $N=3200$\\
\midrule
$2$   & 1 & 1 & 1.04000 & 1.01000 & 1.00250 & 1.00063\\
$4$   & 2 & 1 & 0.99840 & 0.99990 & 0.99999 & 1.00000\\
$6$   & 2 & 2 & 0.93619 & 0.98475 & 0.99623 & 0.99906\\
$12$  & 3 & 3 & 0.81261 & 0.95325 & 0.98833 & 0.99708\\
$30$  & 3 & 6 & 0.62335 & 0.90105 & 0.97506 & 0.99375\\
$210$ & 4 & 24& 0.22866 & 0.74297 & 0.93229 & 0.98288\\
\bottomrule
\end{tabular}
\end{center}

On the ring the budget is spent exactly rather than at most, which turns the
same law into a \textbf{parity selection rule}:
\begin{equation}
n^{\rt{pbc}}_N(\sigma) \;\simeq\;
2P(\sigma)\frac{(N/2)^{\Omega(\sigma)}}{\Omega(\sigma)!}
\quad\text{when } N \equiv \mathrm{sopfr}(\sigma)+\Omega(\sigma) \pmod 2,
\end{equation}
($\mathrm{sopfr}$ is the sum of prime factors with multiplicity, A001414).
The other parity is not merely one degree lower --- usually it is \emph{empty}.
Merging two factors $a,b$ changes $\mathrm{sopfr}+\Omega$ by $(a-1)(b-1)-2$,
which is odd exactly when $a$ and $b$ are both even, so a factorisation of the
opposite parity exists iff $\sigma$ admits one with two even factors:
\begin{equation}
\sigma \text{ occurs at both parities of } N \iff 4 \mid \sigma ,
\end{equation}
and otherwise the wrong parity is empty. (Verified for $\sigma=2\ldots40$ over
$N=3\ldots59$; the sharpening is Tier-2 R-T20 rev.3's.) The effect is not
subtle. At
$N=400$ the ring of $W_{156}$ has \emph{no} sector of size $2$ and none of size
$6$, while at $N=401$ it has $401$ of size $2$; size $4$ occurs $39600$ times at
$N=400$ and $401$ times at $N=401$. For $\sigma=1$ the two frozen states of the
ring defect are present at every $N$ and are added to the formula, reproducing
the exact counts $4$ and $2$.

\subsection{The hyperbolic family: exponentially many}

Here the padding is degenerate --- $K(g)=1$ for $g\in\{0,2,3\}$, at costs $1$, $3$
and $4$ letters --- and $x+x^3+x^4=1$ at $x=1/\varphi$. So the number of ways to
pad is itself exponential and
\begin{equation}
\boxed{\; n_N(\sigma)\;\simeq\;A(\sigma)\,N^{m(\sigma)}\varphi^{N} \;}
\end{equation}
with $m(\sigma)$ the largest number of factors in an ordered factorisation of
$\sigma$ into Fibonacci numbers $\ge2$, and $m(\sigma)=-1$ (no sectors, ever)
when $\sigma$ is not such a product. Measured $n_N(\sigma)\,N^{-m}\varphi^{-N}$
at \rt{obc0} settles to a constant: for $\sigma=1$ ($m=0$) it is $1.3708$ at
$N=100$, $200$ and $400$ alike; for $\sigma=2$ ($m=1$) it moves from
$6.450\times10^{-2}$ to $6.467\times10^{-2}$ over the same window; for $\sigma=8$
($m=3$, $8=2\cdot2\cdot2$ with $2=F(3)$) from $2.04\times10^{-5}$ to
$2.28\times10^{-5}$, the residual drift being the expected $O(m/N)$.

The contrast is the cleanest single statement in this report. \emph{At fixed
size, the number of sectors is polynomial in $N$ for $W_{156}/W_{198}$ and
exponential for $W_{108}/W_{201}$} --- and it is the same fact as
$n_{\rm frozen}$ being linear in one family and $\Theta(\varphi^N)$ in the other,
read at $\sigma=1$.

\begin{figure}[htbp]
\centering
\ifgraphic{../../figures/fig_sector_multiplicity.pdf}{\textwidth}
\caption{\textbf{$n_N(\sigma)$, the number of sectors of size exactly
$\sigma$.}
(a) the measured histogram --- union-find over the $2^N$ flip graph at $N=16$,
\rt{obc0} --- with the renewal formula overlaid as black dashes; they agree
exactly, size by size. The gaps in the $W_{108}$ comb are the non-Fibonacci
integers, which are not sector sizes at any $N$ \emph{on the chain}; the
$W_{156}$ comb has no gaps. On the ring one extra size per $N$ is filled in by
the wall-free sector, the Lucas number $L(N)$.
(b) the same on the ring at $N=16$ and $N=17$: the two combs are largely
\emph{disjoint}, which is the parity selection rule --- $\sigma$ is populated at
leading order only when $N\equiv\mathrm{sopfr}(\sigma)+\Omega(\sigma)\pmod2$.
(c) the parabolic law: $n_N(\sigma)$ against $N$ on log-log, with the closed form
dashed. The slopes are $\Omega(\sigma)+1$, and $\sigma=6$ sits a factor $2$
above $\sigma=4$ at the same slope.
(d) the hyperbolic law: $n_N(\sigma)\varphi^{-N}$ against $N$, straight lines of
slope $m(\sigma)$, with $A(\sigma)N^{m}$ read off at the largest $N$.}
\label{fig:mult}
\end{figure}

\section{The distribution at large $N$}

\subsection{One scalar equation}

By \eqref{eq:gf} the singularity of both generating functions is the root of
$1-B_s$, so
\begin{equation}
\boxed{\;Z_N(s)\;\sim\;\text{const}\cdot\lambda_s^N,\qquad
B_s(1/\lambda_s)=1 \;}
\end{equation}
with
\[
B_s(x) = x\,\mathrm{Li}_{-s}(x) \ \ \text{(parabolic)},\qquad
B_s(x) = x + \sum_{g\ge2}F(g-1)^s x^{g+1}\ \ \text{(hyperbolic)} .
\]
Two things follow immediately. First $\lambda_s$ is \textbf{independent of the
boundary condition}: the boundary enters \eqref{eq:gf} only through the numerator
$A_sC_s$ and the additive term $E_s$, never through the pole. Second the
landmarks are exact: $B_0(x)=x^2/(1-x)$ gives $\lambda_0=\varphi$ and
$B_0(x)=x+x^3/(1-x)$ gives $x^3-x^2+2x-1=0$, i.e.\ $\lambda_0=\rho^2$, while
$B_1$ gives $\lambda_1=2$ in both families.

\begin{center}
\small
\iftable{tab_r30_lambda.tex}
\end{center}

\noindent The minimal polynomials are Tier-2 R-T20's, obtained there from the
$s$-fold tensor transfer matrix $T_s = M_0^{\otimes s}+M_1^{\otimes s}$. Their
dominant roots substituted into $B_s(1/\lambda)$ return $1$ to $30$ digits, so
the two constructions describe the same object; the tensor route was also
reimplemented here and reproduces the $s$-th moment of the exact histogram for
$s=1\ldots4$, all four rules, both boundary conditions.

\subsection{The sizes are log-normal}

A label is a renewal chain and $\ln|\K| = \sum_j \ln K(g_j)$ is a
renewal-reward sum, so the central limit theorem applies to $\ln|\K|$ under the
uniform measure on sectors. With $\tau(s)=\ln\lambda_s$,
\begin{equation}
\boxed{\;|\K| \;=\; \exp\bigl(\alpha_{\rm typ}N + \sigma\sqrt{N}\,\xi\bigr),
\qquad \xi\sim\mathcal{N}(0,1),\qquad
\alpha_{\rm typ}=\tau'(0),\ \ \sigma^2=\tau''(0). \;}
\end{equation}
Beyond the typical value the whole spectrum follows: the number of sectors of
size $e^{\alpha N}$ is $e^{Ng(\alpha)}$ with
$g(\alpha)=\inf_s\bigl(\tau(s)-s\alpha\bigr)$. Its right edge is
$\alpha_{\max}=\lim_{s\to\infty}\tau'(s)$, where $g=0$ and a single sector
survives --- $\ln\varphi$ in the hyperbolic family and $(\ln4)/5$ in the
parabolic one, matching \S3.2 items 1 and 2. Its left edge is $\alpha=0$, where
$g(0)=\ln\lambda_{-\infty}$ is the frozen-sector growth rate --- $\ln\varphi$ and
$0$ respectively, matching item 3. A separate landmark is $\tau'(1)$: the
$s=1$ measure weights each sector by its size, i.e.\ it is the uniform measure on
\emph{states}, so $e^{\tau'(1)N}$ is the size of the sector a randomly chosen
state finds itself in. That $\tau'(1)<\ln2$ strictly is exactly strong
fragmentation in the sense of R21.

\begin{center}
\small
\resizebox{\textwidth}{!}{\iftable{tab_r30_asymptotics.tex}}
\end{center}

\begin{figure}[htbp]
\centering
\ifgraphic{../../figures/fig_sector_sizes.pdf}{\textwidth}
\caption{\textbf{The sector-size distribution.}
(a) the exact distribution on the ring at $N=60$ against the predicted
log-normal $\mathcal{N}(\alpha_{\rm typ}N,\sigma^2N)$, no fitted parameter. The
dotted lines mark $\alpha_{\max}$, which the distribution never reaches at any
appreciable density even though a sector of that size exists.
(b) the large-deviation spectrum $g(\alpha)$, with the three landmarks: the
typical sector ($s=0$, circle), the typical \emph{state's} sector ($s=1$,
square), and the largest sector ($\alpha_{\max}$, triangle, where $g=0$). The
parabolic curve approaches its right edge only logarithmically in $s$, so the
plotted branch stops just short of the triangle.
(c) the finite-size offset $\langle\ln|\K|\rangle-\alpha_{\rm typ}N$: identically
zero on the ring for all four rules, and a rule-dependent constant on the chain
--- \emph{positive} for $W_{201}$.}
\label{fig:sizes}
\end{figure}

\subsection{The two families, side by side}

For $W_{156}/W_{198}$ the three landmarks $\alpha_{\rm typ}=0.198741$,
$\tau'(1)=0.223285$ and $\alpha_{\max}=0.277259$ are close together and
$\sigma^2=0.030461$ is small: the sectors are all of roughly the same size. For
$W_{108}/W_{201}$ they are $0.074455$, $0.208469$ and $0.481212$ with
$\sigma^2=0.076019$: the distribution is far broader and far more unequal. The
crudest way to see it is to compare the mean sector with the sector a random
state lands in --- at $N=16$, \rt{obc0}, the mean $W_{201}$ sector holds $11.2$
states but a random state sits in one holding $200.8$, against $25.4$ and $40.0$
for $W_{156}$. This is the same linear-versus-exponential gap module that
separates the two families in R-T17 and R27, now visible in the census.

\subsection{The boundary: exact on the ring, rule-dependent on the chain}

Since $\lambda_s$ is boundary-independent, so are $\alpha_{\rm typ}$ and
$\sigma^2$; the boundary can only contribute $O(1)$. Writing
$\langle\ln|\K|\rangle = \alpha_{\rm typ}N + c_1$ and
$\mathrm{Var}\,\ln|\K| = \sigma^2 N + c_2$, an exact $O(N^2)$ renewal dynamic
program over the $L/K/R$ decomposition gives
\begin{center}
\small
\begin{tabular}{lrrrr}
\toprule
& $W_{156}$ & $W_{198}$ & $W_{108}$ & $W_{201}$\\
\midrule
$c_1$ (\rt{obc0}) & $-0.246953$ & $-0.246953$ & $-0.100987$ & $+0.245187$\\
$c_2$ (\rt{obc0}) & $+0.199433$ & $+0.199433$ & $-0.148880$ & $+0.170421$\\
$c_1,c_2$ (\rt{pbc}) & $0$ & $0$ & $0$ & $0$\\
\bottomrule
\end{tabular}
\end{center}

\begin{remark}[the offset is not a universal pad cost]
An earlier draft of this analysis measured $c_1=-0.2470$ for $W_{156}$ and
described it as ``the pad cost''. That reading is wrong, and Tier-2 R-T20 rev.2
caught it: there are four values and $W_{201}$'s has the \emph{opposite sign}.
The mechanism is the same one that separates $F(N+2)$ from $F(N)$ in \S3.2 item
1 --- the frozen $0$-pad is a free collar for the word \rt{11} and a block letter
for \rt{00} --- now showing up in the $O(1)$ term of the central limit theorem
rather than in a leading exponent. The values above are reproduced independently
here and by Tier-2 and agree to seven digits. No closed form is known for any of
the three; $-0.24695309$ is not $\ln(25/32)=-0.2468601$.
\end{remark}

\noindent On the ring both offsets vanish identically --- the residuals at
$N=192$ are $\sim10^{-6}$ and drift linearly in $N$, which identifies them as the
truncation error of the substituted $\alpha_{\rm typ}$ and $\sigma^2$ rather than
a genuine constant. So the log-normal law holds on the ring at \emph{finite} $N$
with no correction at all: at $N=24$ the measured mean is $0.198739N$ against
$\alpha_{\rm typ}=0.198741$, and the state-weighted mean is $0.223285N$ on the
nose.

\section{What is proved, what is computed, what is not claimed}

\paragraph{Proved.} Theorem~\ref{thm:label} (the label is the sector, with the
one ring defect) and Proposition~\ref{prop:renewal} (rank one gives the renewal
factorisation). The achievable-size statement in \S3.2 item 6 was falsified in
its first form --- it omitted the wall-free label, which is a trace on the ring
and therefore Lucas rather than Fibonacci --- by Tier-2 R-T20 rev.3, and is
stated here in its corrected form; the routines were already right, and the
regression test that should have caught it was vacuous because its list of
system sizes happened to miss the $N$ at which $L(N)$ is small. Given those, the generating functions \eqref{eq:gf}, the closed
forms for $D_{\max}$, $n_{\rm frozen}$ and $n_{\rm sectors}$, the achievable-size
semigroups, and $\lambda_1=2$ are derivations.

\paragraph{Computed exactly, and cross-checked three ways.} The size multisets
agree between (i) union-find on the $2^N$ flip graph, (ii) the renewal product
enumerated label by label, and (iii) the polynomial-state transfer dynamic
program, for all four rules, both boundary conditions and $N\le16$; and the
moments agree with the independent tensor construction for $s=1\ldots4$.
Tier-2 R-T20 reproduces the same multisets from its own engine, and its eight
minimal polynomials satisfy $B_s(1/\lambda_s)=1$ to $30$ digits.

\paragraph{Asymptotic, with a measured rate of approach.} The log-normal law and
the large-deviation spectrum are $N\to\infty$ statements; the tables above give
the finite-$N$ residuals rather than asserting convergence. The
$n_N(\sigma)$ closed forms are leading-order at fixed $\sigma$ and their
$1/N$ corrections are visible in the ratio table --- $\sigma=210$ is still
$1.7\%$ low at $N=3200$.

\paragraph{Not claimed.} No closed form for the three chain offsets $c_1$, nor
for the residue $A(\sigma)$ in the hyperbolic multiplicity law; both are
available numerically to any precision from the exact routines. Nothing here
extends to the other twelve unitary rules, and nothing here was obtained by
invoking a reflection or a spin flip.

\paragraph{A caveat on cost, corrected.} An earlier note described
\fn{hist} as polynomial in $N$. It is not: it carries one state per distinct
partial matrix product, which is $1.0\times10^4$ at $N=40$, $6.4\times10^5$ at
$N=80$ and $1.3\times10^7$ at $N=120$ (about $5.8$ GB). The number of distinct
\emph{sizes} is polynomial ($540$ at $N=32$); the number of partial products is
not. Above $N\approx100$ use \fn{moment\_Z} for moments, \fn{cumulants} for the
central limit constants and \fn{multiplicity} for $n_N(\sigma)$ --- all three are
polynomial or better.

\section{Reproduction}

\begin{verbatim}
python -m qca_fragmentation.scaling.sector_sizes --bc obc0 --nmax 22
python -m qca_fragmentation.scaling.sector_sizes --bc pbc  --nmax 22
python -m qca_fragmentation.scaling.sector_size_figure
python -m pytest qca_fragmentation/tests/test_sector_sizes.py -q     # 160 tests
\end{verbatim}

\noindent Data in \fn{analytics/sector\_sizes\_\{obc0,pbc\}.json}. The
complementary transfer-matrix treatment, the minimal polynomials of $\lambda_s$,
the $G(B)$ residue formula for $D_{\max}$ and the independent offset table are
Tier-2 report R-T20 rev.2 in the \fn{sector\_analysis\_tier2} repository.

\end{document}
